% P10-4b, 14.09.2026: Begriffs-/Beleganschlüsse und abgegrenzte Prüfaussagen präzisiert.
% ============================================================
% Tabelle 7.3-2: 3x3-Übergangstabelle der Stufenanalyse — v02 (08.09.)
% In Overleaf ablegen als: tables/eval-stufen-uebergang.tex
% Aufruf im Fließtext: \input{\tabledir/eval-stufen-uebergang.tex}
% Label: t:eval-stufen
% Stil: chap6.1-Konvention.
% HERKUNFT: runs/w2-z1-stufenanalyse-8189ea.json (alle 109 A-/B-
% Urteile mit Beleg) + runs/w2-nachreview-konsolidierung.json
% §aBAbgleich (A-Urteile der 5 geänderten IDs 025/062/086/105/107
% unter denselben Auslegungen nachgezogen — alle folgen ihrer
% B-Klasse; Klassenwechsel bleiben exakt G-IE-021/023).
% Lauf: 20260905_164650_8189ea.
% ============================================================

\begin{table}[htbp]
  \centering
  \small
  \begin{tabular}{lrrrr}
    \hline
    & \multicolumn{3}{c}{\textbf{Endbestand B (47 Claims)}} & \\
    \textbf{Kandidaten A (54 Claims)} & full & partial & none & $\Sigma$ A \\
    \hline
    full    & 54 & 1 & 1 & 56 \\
    partial & 0 & 28 & 0 & 28 \\
    none    & 0 & 0 & 25 & 25 \\
    \hline
    $\Sigma$ B & 54 & 29 & 26 & 109 \\
    \hline
  \end{tabular}
  \caption[Abdeckung vor und nach der Kanonisierung im Stresslauf]{Übergang der Abdeckungsklassen aller 109 Referenzaussagen
    zwischen Extraktions-Kandidaten (A) und kanonischem Endbestand
    (B) desselben Stresslaufs, im konsolidierten Auswertungsstand
    (A-Urteile der fünf neu entschiedenen Aussagen unter denselben
    Auslegungen nachgezogen). Nur zwei Aussagen wechseln die Klasse
    (beide Verluste, beide aus derselben Zusammenführung); die
    kleinere Claim-Zahl selbst ist kein Verlustnachweis.}
  \label{t:eval-stufen}
\end{table}
